\documentclass[10pt]{article}
\usepackage[T1]{fontenc}
\usepackage[utf8]{inputenc}
\usepackage[french]{babel}
\usepackage{amsmath,amssymb,mathrsfs,mathpazo}
\usepackage{geometry}
\geometry{a4paper,margin=2cm}
\usepackage{graphicx,array,tabularx,tabularray,colortbl,multirow,multicol,booktabs,enumitem,pifont,ulem,fancybox,diagbox,cancel}
\usepackage{pstricks,pstricks-add,pst-plot,pst-tree,pst-node,pst-text,pst-func,pst-math,pst-bspline,pst-3dplot,pst-eucl,pst-solides3d}
\let\pscalculate\undefined
\usepackage{pst-bezier}
\usepackage{tikz,pgfplots,tkz-tab}
\usetikzlibrary{arrows,arrows.meta,patterns,calc,positioning,trees,decorations.markings,shapes.geometric,shapes.misc,shapes.arrows,intersections,decorations.pathmorphing,decorations.pathreplacing,backgrounds,shadows,babel}
\pgfplotsset{compat=1.18}
\usepgfplotslibrary{fillbetween}
\usepackage[np]{numprint}
\usepackage{listings,makecell,pas-tableur,twoopt}
\let\barre\undefined
\usepackage[pstricks]{tablvar}
\definecolor{darkgreen}{RGB}{0,100,0}
\providecommand{\red}{\color{red}}
\providecommand{\blue}{\color{blue}}
\providecommand{\green}{\color{green}}
\providecommand{\black}{\color{black}}
\providecommand{\white}{\color{white}}
\providecommand{\gray}{\color{gray}}
\providecommand{\cyan}{\color{cyan}}
\providecommand{\magenta}{\color{magenta}}
\providecommand{\yellow}{\color{yellow}}
\providecommand{\violet}{\color{violet}}
\providecommand{\orange}{\color{orange}}
\providecommand{\euro}{\texteuro}
\providecommand{\EUR}{\texteuro}
\providecommand{\ieme}{\textsuperscript{e}}
\providecommand{\degres}{\ensuremath{^\circ}}
\providecommand{\degree}{\ensuremath{^\circ}}
\providecommand{\vv}[1]{\overrightarrow{#1}}
\providecommand{\Coord}[1]{#1}
\providecommand{\decosix}{}
\providecommand{\starredbullet}{\ensuremath{\bullet}}
\pagestyle{empty}
\setlength{\parindent}{0pt}
\usepackage{adjustbox}\DeclareUnicodeCharacter{2236}{:}\DeclareUnicodeCharacter{2190}{\ensuremath{\leftarrow}}
\providecommand{\R}{\mathbb{R}}
\providecommand{\N}{\mathbb{N}}
\providecommand{\Z}{\mathbb{Z}}
\providecommand{\Q}{\mathbb{Q}}
\providecommand{\D}{\mathbb{D}}
\providecommand{\C}{\mathbb{C}}
\providecommand{\e}{{\mathrm{e}}}
\providecommand{\vect}[1]{\overrightarrow{#1}}
\providecommand{\vectt}[1]{\overrightarrow{\mathrm{#1}}}
\providecommand{\vvt}[1]{\overrightarrow{\mathrm{#1}}}
\providecommand{\np}[2][]{{#2~#1}}
\providecommand{\vv}[1]{{\overrightarrow{#1}}}
\renewcommand{\d}{{\mathrm{d}}}
\providecommand{\degres}{{^\circ}}
\providecommand{\degre}{{^\circ}}
\providecommand{\blue}{\color{blue}}
\providecommand{\red}{\color{red}}
\providecommand{\enskip}{\quad}
\providecommand{\llbracket}{\lbrack\!\lbrack}
\providecommand{\rrbracket}{\rbrack\!\rbrack}
\providecommand{\euro}{{\text{€}}}
\providecommand{\no}{{\text{n°}}}
\providecommand{\interval}[2]{{[#1;#2]}}
\providecommand{\intervFO}[2]{{[#1;#2[}}
\providecommand{\dotfill}{\ldots}

\providecommand{\ds}{\displaystyle}
\providecommand{\textding}[1]{#1}
\providecommand{\Oij}{\(\mathrm{O};\vec{\imath},\vec{\jmath}\)}
\providecommand{\Oijk}{\(\mathrm{O};\vec{\imath},\vec{\jmath},\vec{k}\)}
\providecommand{\Ouv}{\(\mathrm{O};\vec{u},\vec{v}\)}
\providecommand{\intervFF}[2]{\ensuremath{[#1;#2]}}
\renewcommand{\euro}{\ifmmode\text{\texteuro}\else\texteuro\fi}
\DeclareUnicodeCharacter{00B0}{\ifmmode^{\circ}\else\textdegree\fi}
\DeclareUnicodeCharacter{20AC}{\ifmmode\text{\texteuro}\else\texteuro\fi}
% Énoncés : Éducation nationale. Transcription et corrigés : APMEP, Denis Vergès et contributeurs. Découpage et classement : CoopMaths, Sébastien Lozano, Éric Elter et contributeurs. Adaptation de la présentation : Maths Explorer.
% CC BY-SA : https://coopmaths.fr/www/about/
% https://coopmaths.fr/alea/static/eam/2026/tex/eam_specifique_2026_06_polynesie_3.tex
\begin{document}
\noindent\textit{APMEP / CoopMaths. CC BY-SA. Présentation adaptée par Maths Explorer.}\par\medskip
Aide au calcul pour cet exercice : $8^{3}= 512 ~;~ 12 \times 8^{2}= 768$

\medskip

Dans une ville, l'évolution d'une épidémie en fonction du temps écoulé $x$, exprimé en semaines depuis le début de l'épidémie, est modélisée par une fonction $f$ définie sur l'intervalle $[0~;~12]$.

Pour tout nombre réel $x$ compris entre 0 et $12,\: f(x)$ désigne le nombre de malades observés à l'instant $x$.

La courbe représentative de $f$, notée $\mathcal{C}_f$, est donnée dans le repère du plan ci-dessous.

\begin{center}
\psset{xunit=0.85cm,yunit=0.027cm,arrowsize=2pt 3}
\begin{pspicture*}(-0.9,-40)(13,290)
\psaxes[linewidth=1.25pt,labelFontSize=\scriptstyle,Dy=40]{->}(0,0)(0,0)(13,290)
\multido{\n=0+1}{14}{\psline[linewidth=0.15pt](\n,0)(\n,280)}
\multido{\n=0+40}{8}{\psline[linewidth=0.15pt](0,\n)(13,\n)}
\psplot[plotpoints=2000,linewidth=1.25pt,linecolor=red]{0}{12.024}{x dup mul 12 mul x 3 exp sub 4 add}
\uput[ul](6.5,240){\red $\mathcal{C}_f$}
\uput[d](10,-15){Temps écoulé (en semaines)}\uput[r](0,285){Nombre de malades}
\end{pspicture*}
\end{center}

\begin{enumerate}
\item Lecture graphique
	\begin{enumerate}
		\item Avec la précision permise par le graphique, lire $f(5)$.

		Interpréter ce résultat dans le contexte de l'exercice.

		\item Déterminer graphiquement la semaine où a eu lieu le pic de l'épidémie.

		Avec la précision permise par le graphique, déterminer le nombre de malades cette semaine-là.
	\end{enumerate}

\item On admet que la fonction $f$ représentée ci-dessus est définie par :
\[f(x)=-x^3+12 x^{2}+4\]

	\begin{enumerate}
		\item Calculer $f(1)$.

		Interpréter ce résultat dans le contexte de l'exercice.

		\item Vérifier que pour tout réel $x$ de l'intervalle $[0~;~12]$,
		\[f'(x)= 3 x(- x + 8).\]
		\item Compléter le tableau donné ci-dessous.
		
		
\begin{center}
\begin{tblr}{|l|p{7cm}|}\hline
\centering$x$ & 0  \hfill12 \\ \hline
Signe de $3 x$ &  \\\hline
Signe de $(-x + 8)$ &  \\\hline
Signe de $f'(x)$ &  \\ \hline
~&~\\
Variations de $f$ &  \\
~&~\\ \hline
\end{tblr}
\end{center}
		\item Déterminer la valeur exacte du nombre de malades la semaine où a eu lieu le pic de l'épidémie.
	\end{enumerate}
\end{enumerate}

\end{document}
